\chapter{Drug Use Screening Tests}
\section*{What Is This Test?} Drug-use screening tests detect selected drugs or metabolites in urine, blood, saliva, hair, or other specimens.
\section*{Alternative Names} Toxicology screen; substance-use screening; drugs-of-abuse testing.
\section*{Principle} Initial immunoassays are qualitative or semiquantitative and may cross-react; confirmatory chromatography with mass spectrometry is more specific.
\section*{Why This Test Is Ordered} Screening may support clinical care, safety assessment, pregnancy care, workplace or legal programs, or treatment monitoring, subject to consent and applicable rules.
\section*{Primary vs Secondary Test} A screen is preliminary; unexpected or consequential results require confirmatory testing and medication review.
\section*{How This Test Is Performed} The specimen is collected with chain-of-custody procedures when required and tested for the requested drug classes and detection window.
\section*{What Preparation Is Needed} Report prescriptions, over-the-counter medicines, supplements, and timing of use. Do not change treatment to alter a result.
\section*{Understanding Results} A negative result does not prove no use, and a positive screen does not prove impairment or misuse. Detection depends on substance, dose, timing, specimen, and cutoff.
\section*{Follow-up Tests} Confirmatory testing, medication reconciliation, clinical assessment, and substance-use evaluation may follow.
\section*{References} \begin{enumerate}\item MedlinePlus. Drug Use Screening Tests.\item Substance Abuse and Mental Health Services Administration. Clinical drug testing guidance.\end{enumerate}
\clearpage\section*{Understanding Results and Follow-up}
The laboratory report should identify the specimen, method, units, reference interval or decision limit, and whether the result is positive, negative, abnormal, or inconclusive. Reference intervals vary with age, sex, pregnancy, specimen type, assay, and laboratory; a result outside a range is not automatically a diagnosis, and a result within a range does not always exclude disease. Discuss unexpected results with the ordering clinician, who can decide whether repeat or confirmatory testing, treatment, or referral is appropriate. Do not change medicines or delay urgent care based on an isolated result.
\section*{Regulatory Status and Authorization Date}This chapter describes a test category, not a specific commercial product. FDA, CDSCO, EU IVDR, and MHRA approval or registration dates therefore cannot be assigned without the assay manufacturer, product name, instrument, and intended use. For laboratory-developed tests, an individual agency approval date may not exist. See the Preface for the interpretation of regulatory status.
\clearpage
